\documentclass[conference]{IEEEtran}
\IEEEoverridecommandlockouts

\usepackage[spanish]{babel}
\usepackage[utf8]{inputenc}
\usepackage[T1]{fontenc}
\usepackage{cite}
\usepackage{amsmath}
\usepackage{graphicx}
\usepackage{url}

\begin{document}

\title{SafeRoute: Plataforma de ruteo peatonal optimizado por seguridad para ciudades del Perú}

\author{
\IEEEauthorblockN{Anngie Salazar Alvarez, Luis Alanya Campos, Pedro Vargas Alfaro, Jharvy Cadillo Tarazona}
\IEEEauthorblockA{Facultad de Ingeniería, Universidad Nacional de Ingeniería\\
Curso de Desarrollo de Software, 2026-I\\
Lima, Perú}
}

\maketitle

\begin{abstract}
La inseguridad ciudadana es uno de los principales problemas de las ciudades del Perú, y las aplicaciones de navegación convencionales optimizan por distancia o tiempo, ignorando el riesgo. Este trabajo presenta SafeRoute, una plataforma web que calcula rutas peatonales minimizando la exposición a zonas inseguras. El sistema combina un modelo de riesgo de tres capas ---criminalidad distrital, contexto urbano de OpenStreetMap y reportes ciudadanos verificados--- sobre el grafo de calles, y aplica el algoritmo de Dijkstra ponderado por riesgo. Se implementó con una arquitectura de servicios desacoplados (FastAPI, PostgreSQL/PostGIS, React) y un esquema de pre-cálculo que garantiza respuestas menores a tres segundos. Se implementaron 16 de 19 requisitos funcionales, incluyendo cuentas de usuario, un módulo de reportes en tres modalidades y moderación de evidencia documental.
\end{abstract}

\begin{IEEEkeywords}
ruteo seguro, ciudad inteligente, seguridad ciudadana, sistemas de información geográfica, Dijkstra, OpenStreetMap.
\end{IEEEkeywords}

\section{Introducción}
Varias ciudades peruanas presentan altos índices de criminalidad. El peatón carece de herramientas que le permitan elegir rutas más seguras: aplicaciones como Google Maps o Waze optimizan por distancia y tiempo, no por riesgo. SafeRoute aborda este vacío calculando y visualizando rutas rankeadas por un índice de seguridad, con una ciudad de cobertura configurable; el despliegue de referencia utiliza Lima Metropolitana.

\section{Contexto y fuentes de datos}
Las fuentes oficiales (SIDPOL y el Observatorio del Ministerio del Interior) publican criminalidad agregada por distrito, sin coordenadas exactas ni granularidad temporal. Para compensar la falta de detalle a nivel de calle, SafeRoute añade dos capas: factores de contexto urbano extraídos de OpenStreetMap (comisarías, cámaras, iluminación, densidad de puntos de interés) y reportes ciudadanos verificados georreferenciados a nivel de punto.

\section{Requisitos}
El sistema fue especificado bajo el estándar ISO/IEC/IEEE 29148:2018. Entre los requisitos funcionales destacan: consultar y visualizar la ruta segura, geocodificación de direcciones, ingesta de datos, cálculo del score compuesto y ruteo, gestión de cuentas con JWT, reportes en tres modos, moderación e integración de reportes verificados al score. Al cierre del proyecto se implementaron 16 de 19 requisitos funcionales por completo; los tres parciales corresponden al cálculo programado, la comparación visual de rutas y la diferenciación de reportes en el mapa.

\section{Arquitectura}
SafeRoute emplea servicios desacoplados con la base de datos como único punto de integración: un servicio de ruteo, un servicio de reportes y un cargador de datos. El pre-cálculo desacopla la consulta de las fuentes externas: el cargador ingesta y calcula por lotes, y la API únicamente lee los scores ya almacenados. Cada servicio sigue una arquitectura en capas (API, servicios, repositorios/DAO tras interfaces y base de datos) sobre PostgreSQL con la extensión PostGIS.

Se aplicaron patrones de diseño: Repository y DAO para el acceso a datos, Factory para la creación de repositorios, un contenedor de inyección de dependencias, Strategy para intercambiar fuentes de datos tras interfaces, Template Method para la ingesta y Facade en la capa de servicios, todo sustentado en los principios SOLID.

\section{Implementación}
El score compuesto por segmento combina tres capas ponderadas: criminalidad distrital, contexto urbano y reportes verificados. Estos últimos se asocian por proximidad mediante \texttt{ST\_DWithin}, con un peso según el tipo de incidencia y un decaimiento temporal. El ruteo ejecuta Dijkstra sobre un costo definido como
\begin{equation}
\text{costo}(e) = \text{longitud}(e) \cdot (1 + \alpha \cdot \text{riesgo}(e)),
\end{equation}
devolviendo la ruta más segura y la más corta con métricas comparativas. El módulo de reportes ofrece tres modos (anónimo, con cuenta y con documentos), con autenticación JWT, hashing de contraseñas, rate limiting, validación de archivos (tipo, tamaño y hash SHA-256) y un flujo de moderación por rol. El frontend, construido con React y Leaflet, renderiza rutas coloreadas por riesgo, el panel de score y los reportes como marcadores diferenciados.

\section{Validación}
El backend cuenta con pruebas unitarias del cálculo de scores y del ruteo con grafos sintéticos, pruebas de integración de los reportes y la moderación, y pruebas de extremo a extremo del frontend con Playwright. Se documentó el cumplimiento del OWASP Top 10. Gracias al pre-cálculo y a los índices espaciales GiST, la API de ruteo cumple el objetivo de responder en menos de tres segundos.

\section{Conclusiones y trabajo futuro}
SafeRoute demuestra que es viable un ruteo peatonal consciente del riesgo integrando datos abiertos y aportes ciudadanos, con presupuesto cero. Como trabajo futuro se plantea automatizar la ejecución programada del cargador de datos, mostrar en la interfaz la comparación entre la ruta segura y la corta, enriquecer la diferenciación de reportes en el mapa e incorporar, de publicarse, datos georreferenciados oficiales como capa adicional.

\begin{thebibliography}{1}
\bibitem{iso29148} ISO/IEC/IEEE 29148:2018, \emph{Systems and software engineering --- Requirements engineering}.
\bibitem{sidpol} Plataforma Nacional de Datos Abiertos, \emph{SIDPOL --- Denuncias Policiales}.
\bibitem{mininter} Observatorio Nacional de Seguridad Ciudadana, Ministerio del Interior del Perú.
\bibitem{osm} OpenStreetMap y Overpass API. \url{https://www.openstreetmap.org}
\bibitem{fielding} R. T. Fielding, ``Architectural Styles and the Design of Network-based Software Architectures,'' Cap. 5: REST, 2000.
\bibitem{ley29733} Ley N.º 29733, \emph{Ley de Protección de Datos Personales del Perú}.
\end{thebibliography}

\end{document}
